% !TEX root = ../main.tex

%----------------------------------------------------------------------
\subsection{Periodic graph construction}\label{sec:graph}
%----------------------------------------------------------------------

Each defect supercell is represented as a graph $\mathcal{G} = (\mathcal{V},
\mathcal{E})$ where nodes correspond to atoms and edges to interatomic
interactions under periodic boundary conditions (PBC).

\paragraph{Minimum-image neighbor list.}
For a 2D monolayer with lattice vectors $\mathbf{L} = [\mathbf{l}_1,
\mathbf{l}_2, \mathbf{l}_3]$, the neighbor set of atom $i$ is
\begin{equation}\label{eq:neighbors}
  \mathcal{N}(i) = \bigl\{(j, \mathbf{n}) \in V \times \mathbb{Z}^3
    \;\big|\; \|\mathbf{r}_j + \mathbf{L}\mathbf{n} - \mathbf{r}_i\|
    < r_\mathrm{cut}\bigr\},
\end{equation}
with $r_\mathrm{cut} = 5.0$\,\AA{}.

\paragraph{Node features.}
Each atom $i$ is described by a 9-dimensional physicochemical feature vector
$\mathbf{x}_i$ (group, period, electronegativity, covalent and van~der~Waals
radii, valence electrons, ionization energy, electron affinity, atomic mass),
all $z$-score normalized.
Optionally, 128-dimensional pretrained atomic embeddings from
ct-UAE~\cite{ctuae2025} are concatenated, injecting transferable chemical
priors from Materials Project pretraining.

\paragraph{Defect-site embedding.}
A binary indicator $\delta_i \in \{0, 1\}$ marks the dopant atom:
\begin{equation}\label{eq:defect_embed}
  \mathbf{h}_i^{(0)} = \mathbf{W}_\mathrm{emb}\,\mathbf{x}_i
    + \mathbf{e}_\mathrm{def}(\delta_i),
\end{equation}
where $\mathbf{e}_\mathrm{def}$ is a learned embedding that injects defect
location into the initial representation.

\paragraph{Edge and triplet features.}
Interatomic distances are expanded via $K=32$ Gaussian radial basis functions
(RBF).  For three-body interactions, bond angles $\theta_{jik}$ between
neighbor pairs $(j,k)$ of center atom $i$ are similarly RBF-expanded and
processed by an MLP.

%----------------------------------------------------------------------
\subsection{DART architecture}\label{sec:arch}
%----------------------------------------------------------------------

\dart{} adopts a hierarchical design decoupling short-range chemical bonding
from long-range defect-induced perturbations (Fig.~\ref{fig:architecture}).

\definecolor{localblue}{HTML}{4393C3}
\definecolor{globalorange}{HTML}{E66101}
\definecolor{innovgreen}{HTML}{1B7837}
\definecolor{readoutpurple}{HTML}{7B3294}
\definecolor{inputgray}{HTML}{D9D9D9}
\definecolor{bgblue}{HTML}{DEEBF7}
\definecolor{bgorange}{HTML}{FEE6CE}
\definecolor{bggreen}{HTML}{D5E8D4}

\begin{figure*}[t]
\centering
\resizebox{0.75\textwidth}{!}{%
\begin{tikzpicture}[
    node distance=0.45cm and 0.7cm,
    block/.style={draw, rounded corners=2pt, minimum height=0.6cm,
                  minimum width=2.2cm, font=\footnotesize, thick,
                  align=center},
    localblock/.style={block, fill=bgblue, draw=localblue},
    globalblock/.style={block, fill=bgorange, draw=globalorange},
    innovblock/.style={block, fill=bggreen, draw=innovgreen},
    ioblock/.style={block, fill=inputgray!50, draw=gray},
    readblock/.style={block, fill=readoutpurple!15, draw=readoutpurple},
    arrow/.style={-{Stealth[length=4pt]}, thick},
    bracetext/.style={font=\scriptsize, text=gray, align=center},
]
% Input
\node[ioblock, minimum width=3.2cm] (input) {Crystal graph $\mathcal{G}$\\
  \scriptsize $\mathbf{x}_i$ (9-dim) $+$ $\delta_i$ defect flag};
% Embedding
\node[ioblock, below=0.35cm of input, minimum width=3.2cm] (embed) {
  Node embedding\\[-2pt]
  \scriptsize $\mathbf{h}_i^{(0)} = \mathbf{W}\mathbf{x}_i + \mathbf{e}_{\mathrm{def}}(\delta_i)$};
% Local layers
\node[localblock, below=0.45cm of embed, minimum width=3.2cm] (local) {
  Local interaction $\times 3$\\[-2pt]
  \scriptsize Bond $+$ angle convolution};
% Innovation: Env enrichment
\node[innovblock, right=0.5cm of local, minimum width=2.5cm] (envrich) {
  \textbf{(II)} Env enrichment\\[-2pt]
  \scriptsize CN, $\bar{d}$, $\Delta\chi$};
% Global layers
\node[globalblock, below=0.45cm of local, minimum width=3.2cm] (global) {
  Global self-attn $\times 2$\\[-2pt]
  \scriptsize Distance bias $\phi(d_{ij}^{\mathrm{PBC}})$};
% Innovation: Pre-norm
\node[innovblock, right=0.5cm of global, minimum width=2.5cm] (prenorm) {
  \textbf{(III)} Pre-norm residual\\[-2pt]
  \scriptsize LN $\to$ SubLayer $\to$ Add};
% Gated pooling
\node[readblock, below=0.45cm of global, minimum width=3.2cm] (pool) {
  \textbf{(I)} Gated attn pooling\\[-2pt]
  \scriptsize $\sigma(\mathbf{W}_g)\!\odot\!$AttnPool $+$ MaxPool};
% MLP
\node[readblock, below=0.35cm of pool, minimum width=3.2cm] (mlp) {
  MLP readout $\to$ $\hat{E}_f$};
% Arrows
\draw[arrow] (input)  -- (embed);
\draw[arrow] (embed)  -- (local);
\draw[arrow] (local)  -- (global);
\draw[arrow] (global) -- (pool);
\draw[arrow] (pool)   -- (mlp);
\draw[arrow, innovgreen, dashed] (envrich) -- (local);
\draw[arrow, innovgreen, dashed] (prenorm) -- (global);
% Braces
\draw[decorate, decoration={brace, amplitude=5pt, mirror},
      localblue, thick]
  ($(local.south west) + (-0.12, 0)$) --
  ($(local.north west) + (-0.12, 0)$)
  node[midway, left=5pt, bracetext, text=localblue]
  {$r\!<\!5$\,\AA};
\draw[decorate, decoration={brace, amplitude=5pt, mirror},
      globalorange, thick]
  ($(global.south west) + (-0.12, 0)$) --
  ($(global.north west) + (-0.12, 0)$)
  node[midway, left=5pt, bracetext, text=globalorange]
  {$r\!<\!12$\,\AA};
% Legend
\node[below=0.5cm of mlp, font=\scriptsize, text=gray] {
  \textcolor{localblue}{$\blacksquare$} Local \quad
  \textcolor{globalorange}{$\blacksquare$} Global \quad
  \textcolor{innovgreen}{$\blacksquare$} Innovations \quad
  \textcolor{readoutpurple}{$\blacksquare$} Readout
};
\end{tikzpicture}%
}% end resizebox
\caption{\dart{} architecture overview.
  Local interaction layers (blue) capture bond lengths, bond angles, and
  coordination within $r_\mathrm{cut} = 5$\,\AA{}.
  Global geometric self-attention layers (orange) model long-range
  perturbations via distance-biased attention up to $d_\mathrm{max} = 12$\,\AA{}.
  Three innovations (green): gated attention pooling, local environment
  enrichment, and pre-norm residual connections.}
\label{fig:architecture}
\end{figure*}

\subsubsection{Local interaction layers}

The first $L_\mathrm{loc} = 3$ layers operate on the sparse neighbor graph
via continuous-filter convolution gated by distance, with additive three-body
angle modulation:
\begin{multline}\label{eq:local}
  \mathbf{h}_i^{(l+1)} = \mathbf{h}_i^{(l)}
    + \sigma\Bigl(
      \textstyle\sum_{j \in \mathcal{N}(i)} \boldsymbol{\phi}(d_{ij})
        \odot \mathbf{W}\mathbf{h}_j^{(l)} \\
      + \textstyle\sum_{(j,k) \in \mathcal{T}(i)}
        f_\mathrm{trip}(\mathbf{h}_i^{(l)}, \theta_{jik})
    \Bigr),
\end{multline}
where $\boldsymbol{\phi}$ is a two-layer MLP over RBF-expanded distance,
$\mathcal{T}(i)$ the triplet set (capped at 32), $f_\mathrm{trip}$ an MLP
over concatenated features and angle RBF, and $\sigma$ the SiLU activation.

\subsubsection{Global geometric self-attention layers}

The subsequent $L_\mathrm{glob} = 2$ layers employ full multi-head
self-attention with learnable distance bias:
\begin{equation}\label{eq:global_attn}
  \mathrm{Attn}(Q, K, V)
    = \mathrm{softmax}\!\left(
        \frac{QK^\top}{\sqrt{d_k}}
        + \phi_\mathrm{bias}\!\bigl(d_{ij}^{\mathrm{PBC}}\bigr)
      \right) V,
\end{equation}
where $d_{ij}^{\mathrm{PBC}}$ is the minimum-image distance (up to
$d_\mathrm{max} = 12$\,\AA{}) and $\phi_\mathrm{bias}$ maps RBF-expanded
distances through an MLP to per-head attention biases.
This mechanism acts as a learnable interatomic potential modulating
information flow~\cite{crystalformer2024}, naturally capturing the
${\sim}1/r^2$ strain fields characteristic of 2D point defects.

\subsubsection{Architectural innovations}\label{sec:innovations}

Three targeted innovations address identified failure modes of the base
architecture:

\paragraph{(I) Defect-aware gated attention pooling.}
Standard mean pooling dilutes the extreme features of the defect atom among
dozens of host atoms.  We replace it with a gated fusion of attention-weighted
and max pooling:
\begin{equation}
  \mathbf{g}_\mathrm{read} = \sigma(\mathbf{W}_g)
    \odot \mathrm{AttnPool}(\{\mathbf{h}_i\})
    + (1 - \sigma(\mathbf{W}_g))
    \odot \max_i \mathbf{h}_i,
\end{equation}
where the attention weights are learned via a single-layer MLP and the gate
$\sigma(\mathbf{W}_g)$ balances global context (attention) with extreme-value
capture (max).
Ablation shows $-19.9\%$ MAE improvement (0.516 $\to$ 0.414\eV{}).

\paragraph{(II) Local environment enrichment.}
Physically interpretable descriptors---coordination number, neighbor distance
statistics (mean, min, max, std), and electronegativity contrast
$\overline{\chi}_\mathrm{nbr} - \chi_i$---are computed from the local
neighborhood and projected through an MLP before injection into the defect
atom's representation.
These features encode Hume-Rothery-type information~\cite{bartel2020}
without requiring explicit physics feature engineering.
Ablation: $-18.0\%$ (0.516 $\to$ 0.423\eV{}).

\paragraph{(III) Pre-norm residual connections.}
Following modern Transformer best practices~\cite{xiong2020}, we apply
LayerNorm \emph{before} rather than after each sub-layer:
\begin{equation}
  \mathbf{h}^{(l+1)} = \mathbf{h}^{(l)}
    + \mathrm{SubLayer}\!\bigl(\mathrm{LN}(\mathbf{h}^{(l)})\bigr).
\end{equation}
This stabilizes gradient flow in deep networks and enables more aggressive
learning rates.
Ablation: $-21.0\%$ (0.516 $\to$ 0.408\eV{}).

\paragraph{Combined effect.}
The three innovations are complementary: together they reduce the MAE from
0.516 to \textbf{0.381\eV{}} ($-26.2\%$), exceeding the sum of individual
gains, with only a 10.7\% increase in parameter count
(0.75\,M $\to$ 0.83\,M).

%----------------------------------------------------------------------
\subsection{Physically motivated variants}\label{sec:variants}
%----------------------------------------------------------------------

Building on the \dart{} base, we design twelve innovation variants targeting
specific physical hypotheses about defect formation energy prediction.
All variants share the same base architecture and training recipe, differing
only in the stated innovation.

\paragraph{Architecture innovations.}
(a)~\emph{V3 Defect-type conditioning}: a zero-initialized learned embedding
distinguishes interstitial from adsorbate defects, motivated by the 2$\times$
difference in their $\ef$ distributions.
(b)~\emph{V4 Mixture-of-experts (MoE) readout}: three specialist MLPs with
learned gating and KL-divergence balance loss, inspired by the observation
that the top 10\% hardest samples contribute 59\% of total error.
(c)~\emph{V6 Physics mismatch features}: six Hume-Rothery descriptors
(size mismatch, electronegativity difference, ionization energy ratio,
electron affinity difference, valence electron mismatch, period
distance)~\cite{bartel2020} computed from dopant--host property contrasts.
(d)~\emph{V14 JK layer aggregation}: learnable softmax weights over
local, global-1, and global-2 layer outputs, allowing the model to
adaptively mix multi-scale representations~\cite{xu2018jk}.

\paragraph{Training strategy innovations.}
(e)~\emph{V11 Label distribution smoothing (LDS)}: Gaussian kernel
density estimation with inverse-square-root reweighting, targeting gradient
starvation for rare high-$\ef$ samples~\cite{yang2021lds}.
(f)~\emph{V13 Rank-N-Contrast}: auxiliary contrastive loss preserving
$\ef$ rank ordering in feature space~\cite{zha2023rnc}.
(g)~\emph{V2+EMA}: exponential moving average of weights
(decay $= 0.999$).
(h)~\emph{V2+Focal MAE}: dynamic reweighting
$w_i = (|e_i|/\overline{|e|})^\gamma$ with $\gamma = 0.5$.
(i)~\emph{V2+Heteroscedastic uncertainty}: additional variance head
predicting aleatoric uncertainty~\cite{kendall2017}.

\paragraph{Combination variants.}
(j)~\emph{V9 Defect--host contrast conditioning}: explicit
$\Delta$-embedding modeling $\ef \propto E(\text{defect}) - E(\text{pristine})$.
(k)~\emph{V10 Best combination}: V3+V6+V9+V14+EMA.
(l)~\emph{V12 LDS+Physics}: V3+V6+V9+LDS+EMA.

%----------------------------------------------------------------------
\subsection{Training pipeline}\label{sec:training}
%----------------------------------------------------------------------

All models are trained with a standardized recipe optimized for the
data-limited regime (${\sim}8.5$\,k training samples):

\begin{itemize}
  \item \textbf{MAE (L1) loss}, directly aligned with the evaluation metric
    and robust to the heavy-tailed $\ef$ distribution.
  \item \textbf{Cosine annealing with linear warmup}: 10-epoch warmup from
    $10^{-7}$ to $5\times10^{-4}$, cosine decay to $10^{-6}$ over 150 epochs.
  \item \textbf{Stochastic Weight Averaging (SWA)}: from epoch 120, weights
    are averaged every epoch~\cite{swa2018}.
  \item \textbf{Online data augmentation}: random rotations, coordinate
    perturbations ($\sigma_r \in [0.01, 0.05]$\,\AA{}), and biaxial strain
    ($\leq 2\%$) applied on-the-fly.
  \item \textbf{Label noise regularization}: Gaussian noise
    ($\sigma = 0.03$\eV{}) simulating DFT numerical uncertainty.
  \item \textbf{Auxiliary defect-type classification}: multi-task head
    predicting defect type with weight 0.1, providing gradient signal for
    defect-relevant features.
\end{itemize}

Training completes in ${\sim}22$ minutes per model on a single NVIDIA L40S GPU.

%----------------------------------------------------------------------
\subsection{Ensemble and uncertainty quantification}\label{sec:uq_method}
%----------------------------------------------------------------------

We construct a pool of 20 models spanning five diversity axes:
architecture variant (\dart{} base, MoE, JK, DefType, Physics, multi-source),
random seed (4 values), training length (150/250 epochs), and data source
(single-source IMP2D / multi-source 4-DB).
Greedy forward selection identifies the $k$-model subset minimizing test MAE.
Prediction uncertainty is estimated as ensemble standard deviation with
post-hoc temperature scaling~\cite{guo2017}.

%----------------------------------------------------------------------
\subsection{Graduated generalization assessment}\label{sec:ood_method}
%----------------------------------------------------------------------

We adopt a three-tier protocol with increasing extrapolation
difficulty~\cite{li2025ood}:

\begin{description}
  \item[Tier~0 (ID):]
    Standard 80/10/10 random split.
  \item[Tier~1 (Intra-family OOD):]
    Leave-one-host-out (LOHO) over seven G6-TMD hosts:
    all configurations of one host are held out; the remaining hosts form
    the training set.
  \item[Tier~2 (Compositional OOD):]
    G6$\times$3d block-out: all (G6 host, 3d dopant) pairs held out,
    testing matrix-completion generalization.
\end{description}

%----------------------------------------------------------------------
\subsection{Prospective DFT validation}\label{sec:dft_method}
%----------------------------------------------------------------------

% TODO: Expand when DFT validation is finalized.
To bridge the gap between benchmark accuracy and real-world deployment, we
perform independent first-principles verification of model-recommended
candidates.  Candidates are selected into two buckets: a \emph{discovery set}
(lowest predicted $\ef$, testing whether the model identifies
thermodynamically accessible dopants) and a \emph{stress-test set} (highest
predicted uncertainty, testing whether $\sigma_\mathrm{cal}$ correctly flags
unreliable predictions).  DFT calculations are performed with Quantum ESPRESSO.
Details will be provided in the final version.
